\documentclass[11pt,a4paper,oneside,openany]{book}
\usepackage{biblioflow-ista}

\renewcommand{\titolodocumento}{Documento Change Requests}
\renewcommand{\sottotitolodocumento}{Change Requests}

\begin{document}
\frontmatter
\copertina
\tableofcontents
\listoftables

\mainmatter
\chapter{Introduzione}

\section{Introduzione}

BiblioFlow è un'applicazione web progressiva (PWA) per la biblioteca di ateneo:
permette agli studenti di prenotare un posto studio, di fare il check-in
all'arrivo, di consultare il catalogo e chiedere libri in prestito, e al
personale della biblioteca di gestire posti, prenotazioni, prestiti e utenti.
L'obiettivo principale del sistema è rendere prevedibile l'accesso ai posti
studio, in particolare per gli studenti pendolari e per chi ha esigenze di
accessibilità, e dare al personale una visione aggiornata di occupazione e
anomalie.

Attualmente, l'architettura di BiblioFlow si articola in tre componenti
principali:
\begin{itemize}
  \item un'\textbf{area studente}, installabile come PWA, per la prenotazione
        guidata dei posti sulla mappa della biblioteca, il check-in con QR code,
        il catalogo, i prestiti e le notifiche;
  \item un'\textbf{area staff} per bibliotecari e amministratori: gestione di
        prenotazioni, posti, prestiti, richieste e utenti, anomalie, statistiche
        e scanner dei QR code;
  \item un \textbf{backend} di route handler Next.js che accede a PostgreSQL
        tramite Prisma, affiancato da automazioni periodiche (promemoria di
        check-in, avvisi di scadenza dei prestiti, rilascio dei posti non
        occupati).
\end{itemize}

Questo documento definisce le Change Requests (CR) utili all'evoluzione di
BiblioFlow. Le modifiche hanno l'obiettivo di rendere affidabile la
prenotazione dei posti e di non lasciare vuoti i posti che si liberano.

Nello specifico, le richieste di cambiamento sono:
\begin{itemize}
  \item \textbf{Prenotazione sicura e atomica dei posti studio (CR01):} si
        richiede che l'identità di chi prenota sia sempre quella della sessione
        autenticata e che due studenti non possano mai ottenere lo stesso posto
        nella stessa fascia oraria, nemmeno con richieste simultanee. Le regole
        sugli orari devono essere uniche per creazione ed estensione, così da
        eliminare gli intervalli incoerenti accettati oggi.
  \item \textbf{Lista d'attesa con riassegnazione automatica (CR02):} a
        supporto della garanzia introdotta dalla CR01, si richiede che lo
        studente che trova il posto occupato possa mettersi in coda per quel
        posto e quella fascia. Quando il posto si libera, per cancellazione o
        perché il titolare non si presenta, il primo in coda deve ottenerlo
        automaticamente ed esserne avvisato.
\end{itemize}

\chapter{Change Requests}

\section{Change Request n.1}

{\small
\begin{xltabular}{\linewidth}{SY}
  \toprule
  Change Request ID & CR01 \\
  Titolo & Prenotazione sicura e atomica dei posti studio \\
  Data di Identificazione & 27 agosto 2026 \\
  Richiesta da & Mario Celzo, Renato Mancino, Alfonso Maria Ferrara \\
  Categoria & Corrective, Preventive \\
  Priorità & Alta \\
  \midrule
  \multicolumn{2}{l}{\textbf{Descrizione}} \\
  Situazione Iniziale &
  La creazione di una prenotazione prende l'identificativo dell'utente dal
  corpo della richiesta inviato dal browser, e le operazioni sulla singola
  prenotazione (lettura, cancellazione, estensione) non verificano che appartenga
  a chi le richiede. Il controllo di disponibilità e la scrittura sono due
  operazioni separate: due richieste quasi simultanee possono superare entrambe
  il controllo e prenotare lo stesso posto nella stessa fascia, e il database
  non impedisce le sovrapposizioni. Le regole sugli orari sono duplicate tra
  creazione ed estensione: un intervallo con la fine precedente all'inizio viene
  accettato, mentre check-in ed estensione confrontano date e orari salvati in
  rappresentazioni diverse e rifiutano richieste valide. \\
  Situazione desiderata &
  Identità e titolarità ricavate sempre dalla sessione autenticata, sia nelle
  operazioni dello studente sia nelle route dello staff, con risposte coerenti
  (401 senza sessione, 403 per il ruolo non ammesso, 404 per la risorsa
  altrui). Un unico servizio di dominio che valida gli intervalli
  per creazione ed estensione. Creazione atomica: verifica e scrittura nella
  stessa transazione e un vincolo del database che renda impossibili due
  prenotazioni attive sovrapposte sullo stesso posto. In caso di conflitto lo
  studente riceve una risposta comprensibile (409), mai un errore interno. \\
  Impatti o Rischi &
  La modifica tocca lo schema del database (nuovo vincolo da applicare ai dati
  già presenti), l'autenticazione, tutte le route di prenotazione e le route
  dello staff che le leggono o le modificano. Un rischio è l'aumento dei
  rifiuti per conflitto di serializzazione sotto carico, che vanno tradotti in
  messaggi chiari; un altro è la regressione dei flussi di check-in e di
  amministrazione. \\
  Fuori ambito &
  La lista d'attesa e la riassegnazione dei posti liberati (gestite nella
  CR02); catalogo e prestiti dei libri. \\
  \bottomrule
  \caption{Change Request CR01}\label{tab:cr01}\\
\end{xltabular}}


\clearpage
\section{Change Request n.2}

{\small
\begin{xltabular}{\linewidth}{SY}
  \toprule
  Change Request ID & CR02 \\
  Titolo & Lista d'attesa con riassegnazione automatica dei posti \\
  Data di Identificazione & 27 agosto 2026 \\
  Richiesta da & Mario Celzo, Renato Mancino, Alfonso Maria Ferrara \\
  Categoria & Perfective \\
  Priorità & Media \\
  \midrule
  \multicolumn{2}{l}{\textbf{Descrizione}} \\
  Situazione Iniziale &
  Quando il posto scelto è già prenotato nella fascia richiesta, lo studente
  riceve soltanto un errore e deve riprovare più tardi. Se il titolare annulla
  la prenotazione o non si presenta (il no-show viene rilevato
  dall'automazione periodica), il posto torna disponibile senza che nessuno ne
  sia informato e spesso resta vuoto, anche se altri studenti lo avevano
  cercato. \\
  Situazione desiderata &
  Dal conflitto, sulla mappa, sulla griglia per dispositivi mobili o nel
  messaggio di posto non più disponibile, lo studente può entrare in una lista
  d'attesa per quel posto e quella fascia, ordinata per ordine di arrivo e senza
  richieste duplicate. Quando il posto si libera (cancellazione dello studente
  o dello staff, no-show, promozione non confermata) il primo in coda viene
  promosso automaticamente a una prenotazione confermata e riceve una notifica;
  se non conferma con il check-in entro il tempo previsto, la promozione scade
  e il posto passa al successivo. Ogni passaggio resta tracciato e lo staff vede
  l'esito della promozione. \\
  Impatti o Rischi &
  La modifica richiede una nuova entità nel modello dati e l'aggiornamento
  delle automazioni periodiche, delle cancellazioni, delle notifiche, della
  mappa dei posti e delle schermate dello staff. I rischi principali sono le
  doppie promozioni quando il cron e una cancellazione agiscono nello stesso
  momento, le notifiche duplicate e l'alterazione delle statistiche esistenti
  da parte dei nuovi eventi. \\
  Fuori ambito &
  Priorità diverse dall'ordine di arrivo e code generiche per sala o per
  fascia; la CR riguarda solo la coda sul singolo posto e intervallo. Nessuna
  modifica al flusso dei prestiti. \\
  \bottomrule
  \caption{Change Request CR02}\label{tab:cr02}\\
\end{xltabular}}


\end{document}
